\chapter{Free relativistic string}
\section{OCQ for the closed string}
... 

\subsection{Quantization}
In canonical quantization we now have to impose the equal time commutation relation
\begin{equation}
    \comm{X^\mu(\tau,\sigma)}{P^\nu(\tau,\sigma')} = i\eta^{\mu\nu}\delta(\sigma - \sigma')
\end{equation}
on the solutions +EQREF and +EQREF. At the level of oscillators this means 
\begin{equation}
    \begin{cases}
        \comm{\alpha^\mu_n}{\alpha^\nu_m} = \comm{\t{\alpha}^\mu_n}{\t{\alpha}^\nu_m} = n\eta^{\mu\nu}\delta_{n + m,0} \\
        \comm{\alpha^\mu_n}{\t{\alpha}^\nu_m} = 0 \\
        \comm{X_0^\mu}{P^\nu} = i\eta^{\mu\nu}
    \end{cases}.
\end{equation}
If we define 
\begin{equation}
    \alpha^\mu_n = \frac{1}{\sqrt{n}}\alpha^\mu_n,\quad\alpha^{\mu_n\dag} = \frac{1}{\sqrt{n}}\alpha^\mu_{-n},\quad\forall n > 0,
\end{equation}
we then have 
\begin{equation}
    \comm{\alpha^\mu_n}{\alpha^{\nu\dag}_m} = \eta^{\mu\nu}\delta_{n,m}.
\end{equation}
These are therefore infinitely many harmonics oscillators. The vacuum state $\ket{0}$ is defined as 
\begin{equation}
    \begin{cases}
        \alpha_n\ket{0} = 0 \\
        \bra{0}\alpha_{-n} = 0
    \end{cases},\quad\forall n > 0,
\end{equation}
with the reality condition 
\begin{equation}
    \alpha^{\mu\dag}_n = \alpha^\mu_{-n}.
\end{equation}
Taking into account also the center-of-mass Heisenberg algebra $\comm{X_0^\mu}{P^\nu} = i\eta^{\mu\nu}$, the Hilbert space will thus be generated by the vacuum at the definite momentum $\ket{0,P}$, on which oscillators act. The generic basis element will thus be 
\begin{equation}
    \alpha^{\mu_1}_{-n_1}\cdots\alpha^{\mu_q}_{-n_q}\t{\alpha}^{\nu_1}_{-m_1}\cdots\t{\alpha}^{\nu_q}_{-m_q}\ket{0,P}.
\end{equation}
Thanks to the oscillators, such a state can be in any possible reducible Lorentz representation $(\{\mu_1,{\dots},\mu_q\}\{\nu_1,{\dots},\nu_q\})$ of the integer spin. The momentum is also not constrained in any way, in fact this is the most general off-shell state we can write down. To find the physical states, we have to impose the constraints which are the missing equation for the metric, resulting in the vanishing of teh energy-momentum tensor. 

\subsection{Virasoro algebra}
Now we have to implement the missing equation of motion for $h$ (+EQREF?) as constraints
\begin{equation}
    \begin{cases}
        T_{++}(\sigma^+) = -\dfrac{1}{\alpha'}\p_+X^\mu\p_+X_\mu \\[2ex]
        T_{--}(\sigma^-) = -\dfrac{1}{\alpha'}\p_-X^\mu\p_-X_\mu
    \end{cases}.
\end{equation}
Expanding them in the harmonics 
\begin{equation}
    \begin{cases}
        T_{++}(\sigma^+) = -\sum_{n\in\bb{Z}}L_n e^{-in\sigma^+} \\
        T_{--}(\sigma^+) = -\sum_{n\in\bb{Z}}\t{L}_n e^{-in\sigma^-}
    \end{cases},
\end{equation}
we obtain the classical expressions of the so-called \emph{Virasoro operators}
\begin{equation}
    \begin{cases}
        L_n = \dfrac{1}{2}\sum_{k\in\bb{Z}}\alpha^\mu_{n - k}\alpha^\nu_k\eta_{\mu\nu} \\[2ex]
        \t{L}_n = \dfrac{1}{2}\sum_{k\in\bb{Z}}\t{\alpha}^\mu_{n - k}\t{\alpha}^\nu_k\eta_{\mu\nu}
    \end{cases},
\end{equation}
where the zero-mode is defined as 
\begin{equation}
\label{eq: zero-mode oscillator}
    \alpha_0^\mu = \t{\alpha}_0^\mu = \sqrt{\frac{\alpha'}{2}}P^\mu.
\end{equation}
We have now to interpret the above classical expressions as operators. In doing so we have to pay attention that the oscillators $\alpha_n^\mu$ are not always commuting and theerefore the previous classical expressions may be ambiguous. In fact, on a closer look only $L_0$ suffers from a possibile ambiguity: $L_{n \neq 0}$ is 
\begin{equation}
    L_{n \neq 0} = \frac{1}{2}\sum_{k\in\bb{Z}}\alpha_{n - k}\cdot \alpha_k
\end{equation}
and it is non-ambiguous since $\comm{\alpha^\mu_{n - k}}{\alpha^\nu_k} = 0$, while $L_0$ is 
\begin{equation}
    L_0 = \frac{1}{2}\sum_{k\in\bb{Z}}\alpha_{-k}\cdot \alpha_k
\end{equation}
and it is ambiguous since $\comm{\alpha^\mu_{-k}}{\alpha^\nu_k} \neq 0$. A natural possibility to consider is the oscillator normal ordering, where the creation operators are placed on the left and annihilation operators on the right. Thus, we define
\begin{equation}
\label{eq: definition of normal ordered L0}
    \h{L}_0 := \frac{1}{2}\sum_{k\in\bb{Z}}\NO{\alpha_{-k}\cdot\alpha_k} = \frac{1}{2}\alpha_0\cdot\alpha_0 + \sum_{k\ge0}\alpha_{-k}\cdot\alpha_k.
\end{equation}
Since any different ordering convention will result in a constant shift in $L_0$, for the time being we leave this constant, that we call $a$, free and fix it a later stage from consistency. This brings to the physical state constraints
\begin{equation}
\label{eq: physical states constraints}
    \begin{cases}
        (\h{L}_0 - a)\ket{\rem{phys}} = 0 \\
        (\h{\t{L}}_0 - \t{a})\ket{\rem{phys}} = 0
    \end{cases}.
\end{equation}
\begin{notz}
    From now on, if not needed, we will not explicitly write the $\h{\cdot}$ on quatum operators, as well as the normal ordering. Moreover, in the following we will always focus on $L_n$ since the discussion for $\t{L}_n$ is the same.
\end{notz}
Let us now define the \emph{level operator} $\LO$ as
\begin{equation}
    \LO := \sum_{k > 0}\alpha_{-k}\cdot\alpha_k.
\end{equation}
This operator counts the level of oscillator excitation by assigning to each oscillator $\alpha_{-n}$ a level $n$. For a general state
\begin{equation}
\label{eq: level operator on a general state}
    \LO\alpha^{\mu_1}_{-n_1}\cdots\alpha^{\mu_q}_{-n_q}\ket{0,P} = (n_1 + \cdots + n_q)\alpha^{\mu_1}_{-n_1}\cdots\alpha^{\mu_q}_{-n_q}\ket{0,P}.
\end{equation}
\begin{ex} 
    To better understand the formula in equation \eqref{eq: level operator on a general state}, consider $\LO\alpha^\mu_{-4}\ket{0,P}$. The explicit steps to compute it are 
    \begin{equation}
        \begin{split}
            \LO\alpha^\mu_{-4}\ket{0,P} &= (\alpha_{-4}\cdot\alpha_4)\alpha^\mu{-4}\ket{0,P} \\
            &= \eta_{\rho\sigma}\alpha^\rho_{-4}\alpha^\sigma_4\alpha^\mu_{-4}\ket{0,P} \\
            &= \eta_{\rho\sigma}\alpha^\rho_{-4}(4\eta^{\sigma\mu})\ket{0,P} \\
            &= 4\alpha^\mu_{-4}\ket{0,P}
        \end{split}
    \end{equation}
\end{ex}
With this operator we can then write equation \eqref{eq: definition of normal ordered L0}, using also equation \eqref{eq: zero-mode oscillator}, as
\begin{equation}
    L_0 = \frac{\alpha'P^2}{4} + \LO.
\end{equation}
It is useful to rewrite the constraints in \eqref{eq: physical states constraints} by consider the sum and the difference of the two equations. From the difference we easily get 
\begin{equation}
    (L_0 - \t{L}_0)\ket{0,P} = (\LO - \t{\LO})\ket{0,P} = 0,    
\end{equation}
where we assumed that $a = \t{a}$. If we have a generic state $\ket{\LO,\t{\LO}}$ of level $(\LO,\t{\LO})$, previous equation implies the \emph{level-matching condition}
\begin{equation}
\label{eq: level-matching condition}
    \LO = \t{\LO}.
\end{equation}
It simply establishes that a physical state has to contain the same oscillator level on both sides. Therefore, instead of $(\LO, \t{\LO})$, we will simply use $\LO$ to denote the level of the state of interest. On the other hand, from the sum we get 
\begin{equation}
    (L_0 + \t{L}_0)\ket{0,P} = \bigg(\frac{\alpha'P^2}{2} + \LO + \t{\LO}\bigg)\ket{0,P} = 2a\ket{0,P},
\end{equation}
then by using the level-mathcing condition found in \eqref{eq: level-matching condition}, the mass-shell condition $p^2 = -m^2$, and multiplying by $1/2$, we have 
\begin{equation}
    \frac{1}{2}(L_0 + \t{L}_0)\ket{0,P} = \bigg(-\frac{\alpha'm^2}{4} + \LO\bigg)\ket{0,P} = 2a\ket{0,P},
\end{equation}
which implies  
\begin{equation}
    \frac{\alpha'm^2}{4} = \LO - a.
\end{equation}
Therefore, again because of the level-matching condition, the mass-shell of the state $\ket{\LO,\t{\LO}}$ is 
\begin{equation}
\label{eq: mass-shell for state N}
    \alpha'm^2 = 4(\LO - a).
\end{equation}
\begin{obs}
    From this we can see how crucial it is to know the value of $a$; however, the canonical quantization method can not give a value for this unknown constant. Only by following the BRST quantization, the value for $a$ is naturally derived and not imposed. 
\end{obs}
The question now consist in determining if it is possibile to define the constraint $L_n\ket{\rem{phys}} = 0$, $\forall n \in \bb{Z}$. The short answer is no, but to justify it we need to study the $L_n$ operator algebra. We start by observing that 
\begin{equation}
    \begin{split}
        \comm{L_n}{\alpha^\mu_m} &= \frac{1}{2}\sum_{k\in\bb{Z}}\comm{\alpha^\nu_{n - k}\alpha^\rho_k}{\alpha^\mu_m}\eta_{\nu\rho} \\
        &= \frac{1}{2}\sum_{k\in\bb{Z}}[(n - k)\eta^{\nu\mu}\delta_{n - k + m,0}\alpha^\rho_k + k\alpha^\nu_{n - k}\eta^{\rho\mu}\delta_{k + m,0}]\eta_{\nu\rho} \\
        &= \frac{1}{2}(-m\alpha^\mu_{n + m} - m\alpha^\mu_{n + m}) \\
        &= -m\alpha^\mu_{n+m}.
    \end{split}
\end{equation}
Using this easier-to-get result, we can now evaluate the more difficult relation $\comm{L_n}{L_m}$. The full computation is reported in appendix +REF +APPENDIX, but at the end of the day we have 
\begin{equation}
    \comm{L_n}{L_m} = (n - m)L_{n + m}.
\end{equation}
However, this is only true if $n \neq -m$. In fact, if $n = -m$ there is another term that appears in the commutation relation
\begin{equation}
    \comm{L_n}{L_{-n}} = 2nL_0 + \frac{D}{12}n(n^2 - 1),
\end{equation}
but vanishes for $n \neq -m$. Therefore, the full result for $\comm{L_n}{L_m}$ is 
\begin{equation}
\label{eq: Virasoro algebra}
    \comm{L_n}{L_m} = (n - m)L_{n + m} + \frac{D}{12}n(n^2 - 1)\delta_{n + m,0}.
\end{equation}
The commutation relation found in \eqref{eq: Virasoro algebra} defines the \emph{Virasoro algebra}. This is a particular Virasoro algebra since in general the factor $D$ is substituted by a generic factor $c$ called the \emph{central charge}, hence the corresponding additive term  
\begin{equation}
    \frac{c}{12}n(n^2 - 1)\delta_{n + m,0}
\end{equation}
is called the \emph{central extension}. Thus, to be more precise, equation \eqref{eq: Virasoro algebra} defines the Virasoro algebra with central charge $c = D$. 
\begin{obs}
    The structure of a Virasoro algebra is similar to the one of a Lie algebra, but with an infinite number of generators. Also, the central extension is not a concept referred only to Virasoro algebras, for a Lie algebra the central extension is an operator $C$ which commutes with every generator $G_i$ of the Lie algebra, so $\comm{C}{G_i} = 0$, and therefore 
    \begin{equation}
        \comm{G_i}{G_j} = \gamma_{ij}^kG_k + C,
    \end{equation}
    where $\gamma_{ij}^k$ are the structure constants. The set of central extensions are the so-called \emph{Cartan generators}.

    Moreover, the central extension of a Virasoro algebra can be seen, or better is, a 2-dimensional quantum effect, in the sense that if one would compute the classical Poisson brackets between $L_n$ and $L_m$, they would get only the first term in \eqref{eq: Virasoro algebra} and not the central extension. 
\end{obs}
The central extension also vanishes if $n = \{-1,0,1\}$, which means that $L_{-1}$, $L_0$ and $L_1$ do not have the central extension, but they obey the $\alg{sl}{2,\bb{C}}$ algebra defined by the relations
\begin{equation}
    \comm{L_0}{L_1} = -L_1,\quad \comm{L_0}{L_{-1}} = L_{-1},\quad \comm{L_1}{L_{-1}} = 2L_0.
\end{equation}
The $\alg{sl}{2,\bb{C}}$ algebra is a Lie algebra and can be written in a more compact way with 
\begin{equation}
    \comm{L_i}{L_j} = (i - j)L_{i + j},\quad i,j = \{-1,0,1\},
\end{equation}
where the similarity with \eqref{eq: Virasoro algebra} is more explicit. 
\begin{obs}
    This can be useful since the $\alg{sl}{2,\bb{C}}$ algebra can be pu in an almost one-to-one correspondence with the group of global conformal transformations of the Riemann sphere 
    \begin{equation}
        f(z) = \frac{\alpha z + \beta}{\gamma z + \delta},
    \end{equation}
    whose coefficient matrix has determinant
    \begin{equation}
        \det{
            \begin{pmatrix}
                \alpha & \beta \\
                \gamma & \delta
            \end{pmatrix}
        } = 1.
    \end{equation}
    This is no coincidence and it will be understood while discussing the underlying conformal field theory (CFT) on the worldsheet.
\end{obs}
Recall now the zero-momentum vacuum state $\ket{0,P}|_{P = 0} = \ko$. We have that 
\begin{equation}
    \begin{cases}
        P^\mu\ko = 0 \\
        \alpha^\mu_0\ko = 0 \\
        \t{\alpha}^\mu_0\ko = 0
    \end{cases},
\end{equation}
which implies that 
\begin{equation}
    \begin{cases}
        L_1\ko = 0 \\
        L_0\ko = 0 \\
        L_{-1}\ko = 0
    \end{cases}.
\end{equation}
This means that the state $\ko$ is invariant under the action of $\group{SL}{2,\bb{C}}$ and therefore it is called the $\group{SL}{2,\bb{C}}$-\emph{invariant vacuum}. From the point of view of the zero-momentum vacuum $\ko$, the states $\ket{0,P}$ are not vacua, but rather excited states obtained through 
\begin{equation}
    e^{iP\cdot\h{X}_0}\ko = \ket{0,P}.
\end{equation}
We have indeed 
\begin{equation}
    \begin{split}
        \h{P}^\mu e^{iP\cdot\h{X}_0}\ko &= 0 + \comm{\h{P}^\mu}{e^{iP\cdot\h{X}_0}}\ko \\
        &= iP_\nu\comm{\h{P}^\mu}{\h{X}^\nu_0}e^{iP\cdot\h{X}_0}\ko \\
        &= iP_\nu(-i\eta^{\mu\nu})e^{iP\cdot\h{X}_0}\ko \\
        &= P^\mu e^{iP\cdot\h{X}_0}\ko.
    \end{split}
\end{equation}
Therefore, $\ko$ is the true unique vacuum state of the matter Hilbert space. 

We now have all the ingredients to impose the remaining \emph{Virasoro constraints} for $n \neq 0$ in addition to the ones in \eqref{eq: physical states constraints} for $n = 0$. As a first guess, one could think to simply impose 
\begin{equation}
    L_n\ket{\rem{phys}} = 0,\quad \forall n \neq 0.
\end{equation}
However, because of the Virasoro algebra, this definition would result in 
\begin{equation}
    \mybrakett{\rem{phys}}{\comm{L_n}{L_m}}{\rem{phys'}} = 0 = \mybrakett{\rem{phys}}{(n - m)L_{n + m} + \frac{D}{12}(n^3 - n)\delta_{n + m,0}}{\rem{phys'}},
\end{equation}
which would imply that 
\begin{equation}
    \mybraket{\rem{phys}}{\rem{phys'}} = 0,
\end{equation}
that it is obviously not acceptable. Therefore, the next guess could be the weaker constraint 
\begin{equation}
    \mybrakett{\rem{phys}}{L_n}{\rem{phys'}} = 0,\quad\forall n \neq 0.
\end{equation}
that would imply 
\begin{equation}
    L_n\ket{\rem{phys}} = 0,\quad \forall n > 0,
\end{equation}
and the same for $\t{L}_n$. This is indeed the correct way to impose such constraint as it justified by BRST quantization. Thus, we can define the Virasoro constraints as 
\begin{equation}
\label{eq: Virasoro constraints}
    \begin{cases}
        L_n\ket{\rem{phys}} = 0 \\
        \t{L}_n\ket{\rem{phys}} = 0
    \end{cases},
    \begin{cases}
        \bra{\rem{phys}}L_{-n} = 0 \\
        \bra{\rem{phys}}\t{L}_{-n} = 0
    \end{cases},\quad \forall n > 0.
\end{equation}
where we recalled that $L_n^\dag = L_{-n}$.

\subsection{Physical spectrum of the closed string}
To continue without yet introducing BRST quantization, we at least need to assume maybe the most important result coming from it, that is $a = 1$. There is also another important result which is $D = 26$, but for now it is not necessary to consider it. The most general state can be written as 
\begin{equation}
    \ket{\Psi} = G_{\mu_1\cdots\mu_q\nu_1\cdots\nu_q}(P)\alpha^{\mu_1}_{-n_1}\cdots\alpha^{\mu_q}_{-\mu_q}\t{\alpha}^{\nu_1}_{-m_1}\cdots\t{\alpha}^{\nu_q}_{-m_q}\ket{0,P},
\end{equation}
and this is a physical state if and only if:
\begin{enumerate}
    \item the level-matching condition $\LO = \t{\LO}$ is satisfied;
    \item the mass-shell is given by $\alpha'm^2 = 4(\LO - 1)$;
    \item the Virasoro constraints $L_n\ket{\rem{phys}} = 0$ and $\t{L}_n\ket{\rem{phys}} = 0$, $\forall n > 0$, are satisfied.
\end{enumerate}
We can now analyze the first levels. 

\subsubsection{Level \texorpdfstring{$\LO = 0$}{N = 0}}
The fundamental level is for $\LO = 0$, where we have the constraints 
\begin{equation}
    L_0\ket{0,P} = \frac{\alpha'P^2}{4}\ket{0,P} = a\ket{0,P},
\end{equation}
\begin{equation}
    L_n\ket{0,P} = 0,\quad \forall n > 0.
\end{equation}
From the first constraint and using the mass-shell condition $p^2 = -m^2$ we get 
\begin{equation}
    m^2 = -\frac{4}{\alpha'},
\end{equation}
where we have used that $a = 1$. This result is really concerning since it implies the existence of a tachyon in the spectrum, or equivalently means that we are quantizing the theory around an unstable vacuum.
\begin{obs}
    The question whether is there a new stable vacuum in the tachyon potential (just like the stable symmetry-breaking vacuum of the Higgs field of the Standard Model) is as of today an open question. The strategy in the sequel will be to ignore this tachyon problem and use this model to learn as much as we can in this simple bosonic setting. Later on we will perform a similar analysis for the \emph{super-string model} where the tachyon will not be there and where genuine space-time fermions will arise in the spectrum.
\end{obs}

\subsubsection{Level \texorpdfstring{$\LO = 1$}{N = 1}}
The first excited level is for $\LO = 1$ where $\alpha^\mu_{-1}\t{\alpha}^\nu_{-1}\ket{0,P}$ is the only possible state, whose mass is $m^2 = 0$, since $a = 1$. This state has two Lorentz indices of undefined symmetru, therefore it is a reducible representation of the Lorentz group. If we saturate them with a spacetime polarization we get
\begin{equation}
    G_{\mu\nu}(P)\alpha^\mu_{-1}\t{\alpha}^\nu_{-1}\ket{0,P}.
\end{equation}
The tensor $G_{\mu\nu}$ can be split into three Lorentz irreducible representations 
\begin{equation}
\label{eq: polarization tensor irreps}
    G_{\mu\nu} = h_{\mu\nu} + B_{\mu\nu} + \Phi\eta_{\mu\nu},
\end{equation}
where:
\begin{enumerate}
    \item $h_{\mu\nu}$ is a symmetric traceless tensor defined as 
    \begin{equation}
        h_{\mu\nu} = G_{(\mu\nu)} - D^{-1}\eta^{\mu\nu}G^\rho_\rho;
    \end{equation}
    \item $B_{\mu\nu}$ is an anti-symmetric tensor defined as $B_{\mu\nu} = G_{[\mu\nu]}$;
    \item $\Phi$ is a scalar defined as $\Phi = D^{-1}G^\rho_\rho$.
\end{enumerate}
The most generic state can be written integrating over the momentum 
\begin{eqnarray}
    \ket{\Psi} = \int d^DP\,G_{\mu\nu}(P)\alpha^\mu_{-1}\t{\alpha}^\nu_{-1}\ket{0,P}.
\end{eqnarray}
If we now consider the $n \neq 0$ Virasoro constraints in \eqref{eq: Virasoro constraints}, we have that only $L_1$ gives a non-trivial contribution, while $L_{n\ge2}\alpha^\mu_{-1}\t{\alpha}^\nu_{-1}\ket{0,P} = 0$. We find indeed that 
\begin{equation}
    L_1G_{\mu\nu}(P)\alpha^\mu_{-1}\t{\alpha}^\nu_{-1}\ket{0,P} = \frac{1}{2}G_{\mu\nu}(P)\sum_{k\in\bb{Z}}(\alpha_{1 - k}\cdot\alpha_{k})\alpha^\mu_{-1}\t{\alpha}^\nu_{-1}\ket{0,P},
\end{equation}
where the only contributions are for $k = \{0,1\}$, therefore 
\begin{equation}
    \begin{split}
        L_1G_{\mu\nu}(P)\alpha^\mu_{-1}\t{\alpha}^\nu_{-1}\ket{0,P} &= \frac{1}{2}G_{\mu\nu}(P)\sum_{k\in\bb{Z}}(\alpha_{1 - k}\cdot\alpha_{k})\alpha^\mu_{-1}\t{\alpha}^\nu_{-1}\ket{0,P} \\
        &= G_{\mu\nu}(P)\eta_{\rho\sigma}\alpha^\rho_0\alpha^\sigma_1\alpha^\mu_0\t{\alpha}^\nu_{-1}\ket{0,P} \\
        &= G_{\mu\nu}(P)P^\rho\t{\alpha}^\nu_{-1}\ket{0,P}.
    \end{split}
\end{equation}
Thus, $L_1$ and $\t{L}_1$ impose respectively 
\begin{equation}
\label{eq: transversality conditions}
    \begin{cases}
        P_\mu G^{\mu\nu} = 0 \\
        G^{\mu\nu}P_\nu = 0
    \end{cases},
\end{equation}
which are called the \emph{transversality conditions}. Hence, the general state $\ket{\Psi}$ is physical if and only if the mass-shell is $P^2 = 0$ and the transversality conditions \eqref{eq: transversality conditions} are satisfied. In this case we find the following states:
\begin{enumerate}
    \item the graviton $h_{\mu\nu}$ with its physical conditions
    \begin{equation}
        \begin{cases}
            P^2h_{\mu\nu} = 0 \\
            P_\mu h^{\mu\nu} = 0 \\
            h^\mu_\mu = 0
        \end{cases},
    \end{equation}
    which imply (and in fact are equivalent to) the linearized Einstein equations in momentum space 
    \begin{equation}
        P^2h_{\mu\nu} - P^\rho P_{(\mu}h_{\nu)\rho} - P_\mu P_\nu h^\rho_\rho = 0.
    \end{equation}
    \item The \emph{Kalb-Ramond field} $B_{\mu\nu}$ which is a 2-form $B = B_{\mu\nu}dx^\mu \wedge dx^\nu$ and has physical conditions 
    \begin{equation}
        \begin{cases}
            P^2B_{\mu\nu} = 0 \\
            P_\mu B^{\mu\nu} = 0
        \end{cases},
    \end{equation}
    which are equivalent to the EOMs of a 2-form.
    \item The \emph{dilaton} $\Phi$ which is a scalar field that should obey the physical state condition
    \begin{equation}
        P^2\Phi(P) = 0.
    \end{equation}
\end{enumerate}

\subsection{Searching null states}
It is now useful to introduce some definitions. The first one is already known.
\begin{defn}[physical state]
    A physical state $\ket{\rem{phys}}$ is a state that satisfies the physical conditions 
    \begin{equation}
        \begin{cases}
            (L_0 - a)\ket{\rem{phys}} = 0 \\
            L_{n > 0}\ket{\rem{phys}} = 0
        \end{cases}.
    \end{equation}
\end{defn}
However, there are two other types of states.
\begin{defn}[spurious state]
    A spurious state $\ket{\chi}$ is a state generated by negatively-moded Virasoro operators $L_{-n}$ acting on a generic state $\ket{\chi_n}$, namely 
    \begin{equation}
        \ket{\chi} = \sum_{n > 0}L_{-n}\ket{\chi_n},
    \end{equation}
    such that 
    \begin{equation}
        \mybraket{\chi}{\rem{phys}} = 0.
    \end{equation}
\end{defn}
\begin{defn}[null state]
    A null state $\ket{\rem{null}}$ is a state that is both physical and spurious, namely 
    \begin{equation}
        \begin{cases}
            \ket{\rem{null}} = \sum_{n > 0}L_{-n}\ket{\chi_n} \\
            L_{n > 0}\ket{\rem{null}} = 0 \\
            (L_0 - a)\ket{\rem{null}} = 0
        \end{cases},
    \end{equation}
    such that 
    \begin{equation}
        \mybraket{\rem{null}}{\rem{null'}} = 0,\quad \mybraket{\rem{null}}{\rem{phys}} = 0.
    \end{equation}
\end{defn}
This suggests to define an equivalence relation inside the space of physical states, stating the fact thtat the inner product between physical states does not change if we add null states to the physical states, namely
\begin{equation}
    \ket{\rem{phys}} \simeq \ket{\rem{phys}} + \ket{\rem{null}}.
\end{equation}
The freedom of adding null states to physical states will be perceived in the target space as an on-shell \emph{spacetime gauge invariance}. 

At level $\LO = 0$ there are no null states. The first non-trivial one is at level $\LO = 1$. Recalling that $a = 1$, we have $m^2 = 0$ and therefore $P^2 = 0$. Thus, we have physical zero-norm states. If we consider the polarization
\begin{equation}
    G_*^{\mu\nu} = P^\mu\xi^\nu + v^\mu P^\nu,
\end{equation}
with $P \cdot \xi = 0$ and $v \cdot P = 0$, the state is physical and we also have 
\begin{equation}
    ||G_*^{\mu\nu}\alpha^\mu_{-1}\t{\alpha}^\nu_{-1}\ket{0,P}||^2 \simeq P^2 = 0.
\end{equation}
This is a physical state having vanishing norm, but it is spurious too because
\begin{equation}
    \begin{cases}
        P^\mu\xi^\nu\alpha^\mu_{-1}\t{\alpha}^\nu_{-1}\ket{0,P} \propto L_{-1}(\xi \cdot \t{\alpha}_{-1})\ket{0,P} \\
        v^\mu P^\nu \alpha^\mu_{-1}\t{\alpha}^\nu_{-1}\ket{0,P} \propto (v \cdot \alpha_{-1})\t{L}_{-1}\ket{0,P}
    \end{cases}.
\end{equation}
We hence found a null state, ginving an equivalence relation on the physical states' space
\begin{equation}
    G_{\mu\nu} \simeq G_{\mu\nu} + P^\mu\xi^\nu + v^\mu P^\nu,
\end{equation}
with $P \cdot \xi = 0$, $v \cdot P = 0$, and $P^2 = 0$, which consist in a gauge transformation of the polarization tensor $G_{\mu\nu}$. Using the decomposition \eqref{eq: polarization tensor irreps}, we can write 
\begin{equation}
    h_{\mu\nu} \simeq h_{\mu\nu} + P_{(\mu}\theta_{\nu)},
\end{equation}
\begin{equation}
    B_{\mu\nu} \simeq B_{\mu\nu} + P_{[\mu}\Lambda_{\nu]},
\end{equation}
\begin{equation}
    \Phi \simeq \Phi,
\end{equation}
where we defined 
\begin{equation}
    \theta_\mu = \frac{\xi^\mu + v^\mu}{2},\quad \Lambda_\mu = \frac{\xi_\mu - v_\mu}{2}.
\end{equation}
\begin{obs}
    If we had not assumed from the start that $a = 1$, we would have found that the closed string spectrum at the $(1,1)$-level is physically acceptable for $a \le 1$. At the critical value $a = 1$ there are zero-norm states, corresponding to the gauge symmetries. From a purely OCQ perspective $a = \t{a} = 1$ is required to have a properly defined graviton, dilaton and Kalb-Ramond field with their own gauge transformations, associated to null states. While, at this stage, this is still not sufficient to fix $a = 1$ by consistency it is however an indication that the correct choice $a = 1$ is associated with the appearance of null states in the spectrum.
\end{obs}
Having fixed $a = 1$, at higher levels we find massive representations of the Lorentz group with, in general, spin higher than 2. The structure of these massive modes is quite complicated, especially for the closed string. However, the main thing to grasp here is that just as the search for null states at level $\LO = 1$ selected $a = 1$, the search for null states at level $\LO = 2$ will instead select the critical dimension $D = 26$. To see it, consider 
\begin{equation}
    \ket{\chi_2} = (2L_{-2} + 3L_{-1}^2)\t{\theta}_{-2}\ket{0,P},
\end{equation}
where $\t{\theta}_{-2}$ is a generic operator made up of right oscillators $\t{\alpha}$ adding up to $\t{N} = 2$. We assume that $\t{\theta}_{-2}$ has been chosen so that the $\t{L}_{-n}$ Virasoro constraints are already satisfied, namely
\begin{equation}
    \begin{cases}
        (\t{L}_0 - 1)\t{\theta}_{-2}\ket{0,P} = 0 \\
        \t{L}_1\t{\theta}_{-2}\ket{0,P} = 0 \\
        \t{L}_2\t{\theta}_{-2}\ket{0,P} = 0
    \end{cases}.
\end{equation}
Notice that, thanks to the Virasoro algebra \eqref{eq: Virasoro algebra}, it is enough to impose $\t{L}_1 = \t{L}_2 = 0$ to guarantee that $\t{L}_{n > 0} = 0$. While chacking if the spurious state $\ket{\chi_2}$ is also physical, therefore null, explicit calculations show that
\begin{equation}
    \begin{cases}
        (L_0 - 1)\ket{\chi_2} = \bigg(\dfrac{\alpha'P^2}{4} + 1\bigg)\ket{\chi_2} \\[2ex]
        L_1\ket{\chi_2} = 0 \\
        L_2\ket{\chi_2} \propto (D - 26)\t{\theta}_{-2}\ket{0,P}
    \end{cases}.
\end{equation}
Moreover, one can show that $||\ket{\chi}||^2 \propto (26 - D)$, which is zero precisely when $D = 26$, namely when the state is null, as expected.
\begin{obs}
    It can be shown that when $a = 1$ and $D = 26$ there is a proliferation of similar null states which, as previously argued, define equivalence relations between physical states. This is the so-called \emph{critical string theory} and these null states are the hallmark of the \emph{BRST cohomology of the string} which will be indeed consistently realized only for $a = 1$ and $D = 26$.
\end{obs}

\section{OCQ for the open string}
Let us now focus on the open string case. As we have seen in +REF, the EOMs are 
\begin{equation}
    \begin{cases}
        \p_+\p_-X^\mu(\sigma^+,\sigma^-) = 0 \\
        \delta X_\mu\p_\sigma X^\mu|_{\sigma = 0,\pi} = 0.
    \end{cases}
\end{equation}
The open string's endpoints, corresponding to $\sigma = 0$ and $\sigma = \pi$, are characterized by their boundary conditions which, in this case, can be Neumann or Dirichlet conditions; more precisely:
\begin{enumerate}
    \item a \emph{Neumann boundary condition} is defined as $\p_\sigma X^\mu|_{\rem{endpoint}} = 0$, which preserves Poincaré symmetry and it corresponds to the statement that there is no momentum transfer along a certain direction $\mu$ in spacetime;
    \item a \emph{Dirichlet boundary condition} is defined as $\delta X_\mu|_{\rem{endpoint}} = 0$, which breaks interval translational symmetry, since it  corresponds to the statement that the endpoint is fixed at a certain point along the direction $\mu$ in spacetime.
\end{enumerate}
Since there are two types of boundary conditions, the two endpoints of the open string can have both the same one or they can be one Neumann and one Dirichlet. This gives origin to four types of solutions for $X^\mu$ denoted as $X_{NN}$, $X_{DD}$, $X_{ND}$, and $X_{DN}$, where subscripts refer to the boundary conditions; more precisely:
\begin{enumerate}
    \item for $X_{NN}$ we have $\p_\sigma X^\mu|_{\sigma = 0,\pi} = 0$;
    \item for $X_{NN}$ we have $\delta X_\mu|_{\sigma = 0,\pi} = 0$;
    \item for $X_{ND}$ we have $\p_\sigma X^\mu|_{\sigma = 0}$ and $\delta X_\mu|_{\sigma = \pi} = 0$;
    \item for $X_{DN}$ we have $\delta X_\mu|_{\sigma = 0} = 0$ and $\p_\sigma X^\mu|_{\sigma = \pi}$;
\end{enumerate}

\subsection{N-N string}
...

\subsection{D-D string}
...

\subsection{N-D and D-N strings}
...